\documentclass[10pt]{article}
\usepackage[T1]{fontenc}
\usepackage[utf8]{inputenc}
\usepackage[french]{babel}
\usepackage{amsmath,amssymb,mathrsfs,mathpazo}
\usepackage{geometry}
\geometry{a4paper,margin=2cm}
\usepackage{graphicx,array,tabularx,tabularray,colortbl,multirow,multicol,booktabs,enumitem,pifont,ulem,fancybox,diagbox,cancel}
\usepackage{pstricks,pstricks-add,pst-plot,pst-tree,pst-node,pst-text,pst-func,pst-math,pst-bspline,pst-3dplot,pst-eucl,pst-solides3d}
\let\pscalculate\undefined
\usepackage{pst-bezier}
\usepackage{tikz,pgfplots,tkz-tab}
\usetikzlibrary{arrows,arrows.meta,patterns,calc,positioning,trees,decorations.markings,shapes.geometric,shapes.misc,shapes.arrows,intersections,decorations.pathmorphing,decorations.pathreplacing,backgrounds,shadows,babel}
\pgfplotsset{compat=1.18}
\usepgfplotslibrary{fillbetween}
\usepackage[np]{numprint}
\usepackage{listings,makecell,pas-tableur,twoopt}
\let\barre\undefined
\usepackage[pstricks]{tablvar}
\definecolor{darkgreen}{RGB}{0,100,0}
\providecommand{\red}{\color{red}}
\providecommand{\blue}{\color{blue}}
\providecommand{\green}{\color{green}}
\providecommand{\black}{\color{black}}
\providecommand{\white}{\color{white}}
\providecommand{\gray}{\color{gray}}
\providecommand{\cyan}{\color{cyan}}
\providecommand{\magenta}{\color{magenta}}
\providecommand{\yellow}{\color{yellow}}
\providecommand{\violet}{\color{violet}}
\providecommand{\orange}{\color{orange}}
\providecommand{\euro}{\texteuro}
\providecommand{\EUR}{\texteuro}
\providecommand{\ieme}{\textsuperscript{e}}
\providecommand{\degres}{\ensuremath{^\circ}}
\providecommand{\degree}{\ensuremath{^\circ}}
\providecommand{\vv}[1]{\overrightarrow{#1}}
\providecommand{\Coord}[1]{#1}
\providecommand{\decosix}{}
\providecommand{\starredbullet}{\ensuremath{\bullet}}
\pagestyle{empty}
\setlength{\parindent}{0pt}
\usepackage{adjustbox}\DeclareUnicodeCharacter{2236}{:}\DeclareUnicodeCharacter{2190}{\ensuremath{\leftarrow}}
\providecommand{\R}{\mathbb{R}}
\providecommand{\N}{\mathbb{N}}
\providecommand{\Z}{\mathbb{Z}}
\providecommand{\Q}{\mathbb{Q}}
\providecommand{\D}{\mathbb{D}}
\providecommand{\C}{\mathbb{C}}
\providecommand{\e}{{\mathrm{e}}}
\providecommand{\vect}[1]{\overrightarrow{#1}}
\providecommand{\vectt}[1]{\overrightarrow{\mathrm{#1}}}
\providecommand{\vvt}[1]{\overrightarrow{\mathrm{#1}}}
\providecommand{\np}[2][]{{#2~#1}}
\providecommand{\vv}[1]{{\overrightarrow{#1}}}
\renewcommand{\d}{{\mathrm{d}}}
\providecommand{\degres}{{^\circ}}
\providecommand{\degre}{{^\circ}}
\providecommand{\blue}{\color{blue}}
\providecommand{\red}{\color{red}}
\providecommand{\enskip}{\quad}
\providecommand{\llbracket}{\lbrack\!\lbrack}
\providecommand{\rrbracket}{\rbrack\!\rbrack}
\providecommand{\euro}{{\text{€}}}
\providecommand{\no}{{\text{n°}}}
\providecommand{\interval}[2]{{[#1;#2]}}
\providecommand{\intervFO}[2]{{[#1;#2[}}
\providecommand{\dotfill}{\ldots}

\providecommand{\ds}{\displaystyle}
\providecommand{\textding}[1]{#1}
\providecommand{\Oij}{\(\mathrm{O};\vec{\imath},\vec{\jmath}\)}
\providecommand{\Oijk}{\(\mathrm{O};\vec{\imath},\vec{\jmath},\vec{k}\)}
\providecommand{\Ouv}{\(\mathrm{O};\vec{u},\vec{v}\)}
\providecommand{\intervFF}[2]{\ensuremath{[#1;#2]}}
\renewcommand{\euro}{\ifmmode\text{\texteuro}\else\texteuro\fi}
\DeclareUnicodeCharacter{00B0}{\ifmmode^{\circ}\else\textdegree\fi}
\DeclareUnicodeCharacter{20AC}{\ifmmode\text{\texteuro}\else\texteuro\fi}
% Énoncés : Éducation nationale. Transcription et corrigés : APMEP, Denis Vergès et contributeurs. Découpage et classement : CoopMaths, Sébastien Lozano, Éric Elter et contributeurs. Adaptation de la présentation : Maths Explorer.
% CC BY-SA : https://coopmaths.fr/www/about/
% https://coopmaths.fr/alea/static/e3c/2020/tex/e3c_2020_00_sujet19_2.tex
\begin{document}
\noindent\textit{APMEP / CoopMaths. CC BY-SA. Présentation adaptée par Maths Explorer.}\par\medskip
\medskip

On administre à un patient un médicament par injection intraveineuse.

La première injection est de \np[ml]{10}, puis toutes les heures on lui en injecte \np[ml]{1} .

On étudie l’évolution de la quantité de médicament présente dans le sang en prenant le modèle suivant :

\begin{itemize}
\item  on estime que 20\,\% de la quantité de médicament présente dans le sang est éliminée chaque heure ;
 
 \item pour tout entier naturel $n$, on note $U_n$ la quantité de médicament en ml présente dans le sang au bout de $n$ heures.
\end{itemize}
Ainsi, $U_0=10$.
\begin{enumerate}
\item  Justifier que $U_1=9$.
\item Montrer que, pour tout entier naturel $n$, $U_{n+1}=0,8 U_{n}+1$.

On donne ci-dessous la représentation graphique de la suite $\left(U_n\right)$ :

\psset{xunit=0.10cm,yunit=0.7cm,labelFontSize=\scriptstyle,showorigin=false,dotstyle=Mul}
\begin{pspicture}(-5.26,-1)(112,12.6)




\psaxes[linewidth=0.95pt,Dx=10]{->}(0,0)(-5,-0.5)(110,12)
\def\Func{0.8 x exp 5 mul 5 add }
\psplot[plotpoints=100, plotstyle=dots,dotsize=3pt,linecolor=red]{0}{100}{\Func}
\end{pspicture}

\item Conjecturer la limite de la suite $(U_n)$.

On considère l’algorithme suivant :
\begin{center}
\begin{tabular}[]{|l|}
\hline
$U\leftarrow 10$\\
$N\leftarrow 0$\\
Tant que $U > 5,1$ faire\\
\hspace{1.5em}$U\leftarrow 0,8*U+1$\\
\hspace{1.5em}$N\leftarrow N+1$\\
Fin du tant que\\
Afficher $N$ \\
\hline
\end{tabular}
\end{center}
\item À quoi sert cet algorithme ?

\item À l’aide de l’extrait du tableau de valeurs de la suite $\left(U_n\right)$ donné ci-dessous, donner la valeur de $N$ à l’issue de l’exécution de cet algorithme.

\begin{center}
\begin{tabularx}{\linewidth}{|*{8}{>{\centering \arraybackslash} X|}} \hline
$n$&8&9&10&11&12&13&14\\\hline
$U_n$&\np{5,838861}&\np{5,671089}&\np{5,536871}&\np{5,429497}&\np{5,343597}&\np{5,274878}&\np{5,219902}\\\hline
\end{tabularx}

\bigskip

\begin{tabularx}{\linewidth}{|*{8}{>{\centering \arraybackslash} X|}} \hline
$n$&15&16&17&18&19&20&21 \\\hline
$U_n$&\np{5,175922}&\np{5,140737}&\np{5,11259}&\np{5,090072}&\np{5,072058}&\np{5,057646}&\np{5,046117} \\\hline
 \end{tabularx}
 
\bigskip

\begin{tabularx}{\linewidth}{|*{8}{>{\centering \arraybackslash} X|}} \hline
$n$&22&23&24&25&26&27&28\\
 \hline
$U_n$&\np{5,036893}&\np{5,029515}&\np{5,023612}&\np{5,018889}&\np{5,015112}&\np{5,012089}&\np{5,009671}\\\hline
\end{tabularx}
\end{center}
\end{enumerate}

\vspace{0,5cm}

\end{document}
